\slajdid{10-s056}
\begin{frame}[label=10-s056]
\gusttekst\setlength{\abovedisplayskip}{3pt}\setlength{\belowdisplayskip}{3pt}\setlength{\abovedisplayshortskip}{2pt}\setlength{\belowdisplayshortskip}{2pt}\fontsize{9}{9.5}\selectfont
\begin{columns}[c,onlytextwidth]\column{.36\textwidth}\centering\begin{tikzpicture}[x=.34cm,y=.4cm,font=\sitneoznake,>=Latex]\ctikzset{bipoles/length=.6cm}\draw(-.8,.8)--(.5,.8);\draw[->](-.5,.8)--(-1.1,0);\draw(-1.1,0)--(-3.5,0)node[ocirc]{}node[left]{$V_{I1}$};\draw[->](-.25,.8)--(-.7,0);\draw(-.7,0)--(-.7,-.9)--(-2.3,-.9)node[ocirc]{}node[left]{$V_{I2}$};\draw(.25,.8)--(.65,0);\node[below]at(.05,-.05){$T_1$};\draw(0,.8)--(0,2.8)to[R,l_=$R_{B1}$](0,5.2);\node[left]at(-.25,4.6){4k0};\draw(.65,0)--(2.4,0);\draw(2.4,-.5)--(2.4,.5);\draw(2.4,.32)--(3.1,.7)--(3.1,2.8)to[R,l_=$R_{C2}$](3.1,5.2);\node[left]at(2.85,4.6){1k6};\draw[->](2.4,-.32)--(3.1,-.7);\draw(3.1,-.7)--(3.1,-1.1)to[R,l=$R_{E2}$](3.1,-2.7)node[ground]{};\node[left]at(2.85,-1.9){1k0};\node[right]at(2.6,0){$T_2$};\begin{scope}[shift={(5,-1.1)}]\draw(0,-.5)--(0,.5);\draw(-.6,0)--(0,0);\draw(0,.32)--(.7,.7)--(.7,1.2);\draw[->](0,-.32)--(.7,-.7);\draw(.7,-.7)--(.7,-1.6)node[ground]{};\node[right]at(.25,0){$T_3$};\end{scope}\draw(3.1,-1.1)--(4.4,-1.1);\draw(5.7,.1)--(8,.1)node[ocirc]{}node[right]{$V_o$};\draw(5,1.9)--(5,2.9);\draw(3.1,2.4)--(5,2.4);\draw(5,2.72)--(5.7,3.1)--(5.7,3.4)to[R,l_=$R_{C4}$](5.7,5.2);\node[left]at(5.45,4.9){130};\draw[->](5,2.08)--(5.7,1.7);\draw(5.7,1.7)--(5.7,1.05);\draw(5.45,1.05)--(5.95,1.05)--(5.7,.55)--cycle;\draw(5.45,.55)--(5.95,.55);\draw(5.7,.55)--(5.7,.1);\node[right]at(5.4,2.4){$T_4$};\node[right]at(6,.8){$D_4$};\draw(-.3,5.2)--(6.2,5.2)node[right]{$V_{CC}$};\draw(-3,0)--(-3,-1.6);\draw(-3.25,-2.1)--(-2.75,-2.1)--(-3,-1.6)--cycle;\draw(-3.25,-1.6)--(-2.75,-1.6);\draw(-3,-2.1)--(-3,-2.7)node[ground]{};\node[left]at(-3.3,-1.6){$D_1$};\draw(-1.8,-.9)--(-1.8,-1.6);\draw(-2.05,-2.1)--(-1.55,-2.1)--(-1.8,-1.6)--cycle;\draw(-2.05,-1.6)--(-1.55,-1.6);\draw(-1.8,-2.1)--(-1.8,-2.7)node[ground]{};\node[right]at(-1.5,-1.6){$D_2$};\node at(2,-3.7){Standardno TTL kolo};\end{tikzpicture}
\column{.62\textwidth}
1.\quad Diode D1 i D2 na ulazu kola. One su obavezni deo na ulazu u TTL logička kola, i služe da štite ulaz od napona koji su manji od napona mase. Videli smo da ta situacija može da se pojavi kada imamo diferencijatorski efekat u toku signala. Veliki negativan napon na bio kojem ulazu bi rezultovao u velikoj struji kroz tranzistor T1 koji tada radi u direktnom zasićenju i čija je struja direktno proporcionalna ulaznom naponu, Došlo bi prevelike disipacije na ulaznom delu kola pa i do njegovog „pregorevanja“. U toku normalnih režima rada te diode nemaju nikakvu funkciju tako da ćemo ih u našoj analizi „ignorisati“.
\end{columns}\par\smallskip
2.\quad Tranzistor T4, otpornik $R_{C4}$ i dioda D4, Ta konfiguracija naziva se totempol (totem pole) konfiguracijom i cilj je da zamene otporniku u kolektoru izlaznog tranzistora. Videli smo da je bilo „problema“ sa tim otpornikom. Treba da bude mali da bi strujni kapacitet logičke jedinice bio velik, a treba da bude velik da bi na karakteristici imali što veće pojačanje itd. Ideja je da tranzistor T4 radi samo kada bi trebala na izlazu da bude logička jedinica i da obezbeđuje veliki strujni kapacitet. A da je isključen kada bi na izlazu trebala da bude logička nula i da ne smeta strujnom kapacitetu logičke nule. Suštinski to obezbeđuje tranzistor T2 i to je ranije spomenuta njegova druga uloga.

\end{frame}
